\section{Úvod}

\vspace{1.5em}
{ \large 

\noindent\textbf{Rozdelenie práce:}
\vspace{0.75em}

\noindent Práca aj bodové ohodnotenie boli medzi členov tímu rozdelené rovnomerne.

\vspace{0.75em}
\noindent
\begin{tabularx}{\textwidth}{@{} l @{\hspace{1.5em}} >{\raggedright\arraybackslash}X @{}}
  x274855:    &  Syntaktický analyzátor, dynamická pamäť prekladača, spracovanie chýb\\
  xlehots00: &  Lexikálny analyzátor, Analyzátor výrazov, generovanie vstavaných funkcií, dokumentácia\\
  xgreges00: &  Lexikálny analyzátor, Generátor cieľového kódu\\
  xkyclon00: &  Sémantický analýzátor, Tabuľka symbolov\\
\end{tabularx}

\vspace{1.5em}

\noindent\textbf{Popis práce v týme:}

\vspace{0.75em}
\noindent Každý týždeň v nedeľu o 19:00 hodine sa konalo stretnutie členov týmu. Na tomto stretnutí sme si prešli prácou spravenou cez týždeň, ktorá bola prebraná a schválená, následne sa prebralo, čo je ešte potrebné spraviť, prípadne čo má najvyšiu prioritu a rozdelila sa práca na ďalší týždeň.

\vspace{0.75em}
\noindent Schvaľovanie zmien v repozitár: Každá časť prekladača bola vyvíjaná vo vlastnej vetve, pričom pri komplexnejších moduloch sme vytvárali vetvy iteratívne (označené \verb|1.0|, \verb|2.0|, \verb|3.0|). Zmeny sa do vetvy ukladali až potom, čo bolo možné projekt lokálne preložiť. Priebežný vývoj bol ukladaný do vetvy \verb|semnatic-analysis| (neskôr \verb|semantic-analysis-2|), ktorá bola do hlavnej vetvy \verb|main| zlučovaná na nedeľných stretnutiach až po dokončení ucelenej časti projektu, jej schválenáním členmi tímu ako aj úspešním vykonaním všetkých testov.

\vspace{0.75em}
\noindent Testovali sme prostredníctvom troch pripravených testovacích programov (\verb|tiny_prog.scala|, \verb|small_prog.scala|, \verb|medium_prog.scala|), z ktorých každy rozširoval ako dĺžku programu, tak aj komplexnosť príkazov na preloženie prekladačom.
\newpage
\vspace{0.75em}
\noindent Kód sme uniformne formátovali prostredníctvom \verb|.clang-format|, ktorý sme nastavili nasledovne:\\

\vspace{0.75em}
\noindent
\begin{tabular}{l l}
  Odsadenie riadka: & 4 znaky\\
  Maximálna dĺžka riadka: & 130 znakov\\
  Umiestnenie znaku ukazateľa (\verb|*|) a referencie (\verb|&|): & vpravo \\
  Konvencia umiestnenia označenia rozsahu kódu: & Stroustrup \\
  Maximálny počet prázdnych riadkov: & 3
\end{tabular}
\hspace*{0.25cm}
}